\documentclass[11pt]{article}

\usepackage[margin=1in]{geometry}
\usepackage[T1]{fontenc}
\usepackage[utf8]{inputenc}
\usepackage{lmodern}
\usepackage{booktabs}
\usepackage{longtable}
\usepackage{array}
\usepackage{tabularx}
\usepackage{xcolor}
\usepackage{hyperref}

\hypersetup{
  colorlinks=true,
  linkcolor=black,
  urlcolor=blue,
  pdftitle={shravclock Engineering Report},
  pdfauthor={Shravan Sheth}
}

\newcommand{\reportdate}{April 12, 2026}
\newcommand{\repo}{\url{https://github.com/shravansheth/shravclock}}

\title{
  \vspace{-1.5em}
  \textbf{shravclock}\\
  \large Engineering Report and Product-Oriented Build Documentation
}
\author{}
\date{\reportdate}

\begin{document}
\maketitle

\begin{center}
\small Repository: \repo
\end{center}

\section{Document Purpose}
This document provides a comprehensive technical description of the \texttt{shravclock} desk clock firmware project. The report is written as personal build documentation in a polished engineering style, with enough structure to also function as an internal product dossier. It describes the current implementation in detail, records the technical rationale visible in the repository, and separates future work from presently implemented behavior.

Private deployment details such as Wi-Fi credentials, exact installation location, city name, and geographic coordinates are intentionally excluded.

\section{Project Summary}
\texttt{shravclock} is a battery-powered e-paper desk clock built around a Seeed XIAO ESP32-C6 microcontroller, a 4.2-inch monochrome e-ink display, and a DS3231 real-time clock (RTC). The product goal is not merely to display the current time, but to do so in a way that is visually refined, power-conscious, and robust across long unattended operation.

The firmware combines four major capabilities:
\begin{enumerate}
  \item persistent and accurate timekeeping via a hardware RTC,
  \item low-power minute-based wake, render, and deep-sleep operation,
  \item outdoor weather and location-aware presentation over Wi-Fi,
  \item a self-hosted configuration portal that avoids reflashing for normal setup changes.
\end{enumerate}

In practical terms, the device behaves like a quiet ambient object. It wakes shortly before each minute boundary, optionally performs an hourly network update, redraws the e-paper display with either a partial or full refresh, and then returns to deep sleep.

\section{System Scope}
\subsection{Primary Functions}
The current firmware is designed to provide the following user-visible functions:
\begin{itemize}
  \item local time display in 24-hour \texttt{HH:MM} format,
  \item full weekday and calendar date display,
  \item current outdoor weather with iconography and daily low/high temperature,
  \item city name display derived from reverse geocoding,
  \item dark mode display inversion,
  \item browser-based reconfiguration through a captive setup portal.
\end{itemize}

\subsection{Non-Goals in the Current Revision}
The repository indicates several capabilities that are not yet part of the current production firmware:
\begin{itemize}
  \item active battery measurement and on-screen battery indication,
  \item low-battery shutdown or special dead-battery display flow,
  \item adaptive nighttime refresh reduction,
  \item OTA firmware updates,
  \item additional sensors beyond the existing RTC and display stack.
\end{itemize}

\section{Hardware Platform}
\subsection{Core Components}
\begin{longtable}{>{\raggedright\arraybackslash}p{0.27\textwidth} >{\raggedright\arraybackslash}p{0.67\textwidth}}
\toprule
Component & Description \\
\midrule
Microcontroller & Seeed XIAO ESP32-C6 \\
Display & WeAct 4.2-inch SPI e-ink panel, SSD1683 controller, black and white, 400 \(\times\) 300 pixels \\
RTC & DS3231 Mini on I\textsuperscript{2}C \\
User input & Single momentary pushbutton using active-low input with internal pull-up \\
Battery & 3.7 V Li-ion pack, nominally 5000 mAh in the design notes \\
\bottomrule
\end{longtable}

\subsection{Pin Assignment}
\begin{longtable}{>{\raggedright\arraybackslash}p{0.34\textwidth} >{\raggedright\arraybackslash}p{0.2\textwidth}}
\toprule
Signal & GPIO \\
\midrule
SPI MOSI & 18 \\
SPI SCK & 19 \\
E-ink CS & 1 \\
E-ink DC & 2 \\
E-ink RST & 21 \\
E-ink BUSY & 16 \\
RTC SDA & 22 \\
RTC SCL & 23 \\
Button & 17 \\
\bottomrule
\end{longtable}

\subsection{Hardware Observations}
The hardware selection is internally coherent for a low-power always-visible clock:
\begin{itemize}
  \item the DS3231 provides reliable time retention independent of Wi-Fi availability,
  \item the SSD1683 e-paper panel supports partial refresh for low-disturbance minute updates,
  \item the ESP32-C6 provides Wi-Fi connectivity for time sync, weather fetches, and configuration hosting.
\end{itemize}

\textbf{Labeled inference.} The use of a dedicated RTC, together with aggressive deep sleep, indicates that power autonomy and display stability are primary design goals rather than secondary optimizations.

\section{Software Stack}
\subsection{Build Environment}
The repository is organized as a PlatformIO project using the Arduino framework. The principal firmware implementation resides in \texttt{src/main.cpp}.

\begin{longtable}{>{\raggedright\arraybackslash}p{0.3\textwidth} >{\raggedright\arraybackslash}p{0.64\textwidth}}
\toprule
Item & Details \\
\midrule
Build system & PlatformIO \\
Framework & Arduino \\
Target board & \texttt{seeed\_xiao\_esp32c6} \\
ESP32 platform package & \texttt{platform-espressif32} via PlatformIO package URL \\
Partitions & \texttt{min\_spiffs.csv} for maximizing application flash budget \\
\bottomrule
\end{longtable}

\subsection{Key Firmware Dependencies}
\begin{longtable}{>{\raggedright\arraybackslash}p{0.3\textwidth} >{\raggedright\arraybackslash}p{0.64\textwidth}}
\toprule
Library or subsystem & Role \\
\midrule
GxEPD2 & E-paper display driver and rendering support \\
RTClib & DS3231 access and RTC time conversion utilities \\
ArduinoJson v7 & JSON parsing for API responses \\
WiFi / HTTPClient / WiFiClientSecure & network access, HTTP weather fetch, HTTPS geocoding \\
Preferences & non-volatile configuration storage through ESP32 NVS \\
WebServer + DNSServer & captive setup portal implementation \\
\bottomrule
\end{longtable}

\section{Repository Structure}
The repository is intentionally compact. The most important files in the current revision are:

\begin{longtable}{>{\ttfamily\raggedright\arraybackslash}p{0.33\textwidth} >{\raggedright\arraybackslash}p{0.61\textwidth}}
\toprule
Path & Purpose \\
\midrule
src/main.cpp & Main firmware implementation and runtime state machine \\
FIRMWARE\_OVERVIEW.md & High-level architecture and design notes \\
platformio.ini & PlatformIO environment and dependency definitions \\
include/WeatherStationFonts.h & Meteocons bitmap font asset used for weather icons \\
battery-measuring/battery\_code\_backup.cpp & staged battery-measurement feature code not yet integrated \\
README.md & short project summary \\
\bottomrule
\end{longtable}

\section{Functional Specification}
\subsection{Timekeeping}
The displayed time is formatted in 24-hour style and rendered without seconds. Time is stored in UTC on the DS3231 RTC. Local presentation is produced by applying a POSIX timezone string at display time.

This is an important architectural decision:
\begin{itemize}
  \item the RTC remains timezone-agnostic,
  \item configuration changes to timezone rules can be applied without rewriting the conceptual time basis,
  \item daylight saving transitions are handled through timezone conversion logic rather than through local-time RTC storage.
\end{itemize}

On cold boot, the firmware connects to Wi-Fi, performs NTP synchronization, and adjusts the DS3231 using the synchronized UTC epoch.

\subsection{Date Presentation}
The date area contains two centered lines:
\begin{itemize}
  \item full weekday name,
  \item abbreviated month, day of month, and year.
\end{itemize}

The implementation computes text bounds dynamically before drawing, allowing the firmware to center each string without hard-coded assumptions about pixel width.

\subsection{Weather Presentation}
Current weather is sourced from Open-Meteo and shown as:
\begin{itemize}
  \item current temperature in degrees Fahrenheit,
  \item weather icon selected from WMO weather code,
  \item daily low and high temperatures,
  \item city label beneath the weather line.
\end{itemize}

Weather data is not fetched every minute. Instead, the device refreshes weather once per hour near the top of the hour, with retry logic on failure and persistent in-memory caching across deep sleep cycles.

\subsection{Weather Iconography}
The icon system uses a custom renderer for the Meteocons bitmap font format included in the repository. The firmware maps WMO weather codes to a reduced symbolic visual language such as clear, partly cloudy, fog, rain, snow, showers, and thunderstorm.

This approach provides a compact and deterministic icon pipeline without external graphical assets beyond the bundled font table.

\subsection{City Name Resolution}
The city name is obtained through reverse geocoding from Nominatim using the configured latitude and longitude. The firmware prioritizes \texttt{city}, then \texttt{town}, then \texttt{village}, and finally falls back to a generic display name if the more specific fields are absent.

The resulting city label is cached in RTC memory so it can persist across deep sleep cycles without repeated network access.

\subsection{Dark Mode}
Dark mode reverses the display polarity from black-on-white to white-on-black and adds a crescent moon overlay in the top-left area of the screen. The mode is user-configurable through the setup portal and is stored in NVS so that it survives resets and sleep cycles.

\section{Display Layout and Typography}
\subsection{Layout}
The 400 \(\times\) 300 pixel display is arranged as a vertically centered composition:
\begin{itemize}
  \item top emphasis on time,
  \item mid-screen date information,
  \item lower weather block including icon, temperature, low/high values, and city name.
\end{itemize}

The layout is computed programmatically. Bounding boxes are measured at runtime and used to center or align content precisely. This makes the UI resilient to font metrics and reduces the need for manual per-string adjustments.

\subsection{Typography}
The firmware uses a small set of Adafruit GFX fonts with clear role separation:
\begin{itemize}
  \item bold sans serif at double scale for the time,
  \item bold sans serif for weekday,
  \item regular sans serif for date,
  \item bold larger sans serif for current temperature,
  \item smaller sans serif for low/high and city.
\end{itemize}

The degree marker is drawn manually as a circle instead of relying on a font glyph. This gives consistent visual positioning in the weather block.

\textbf{Labeled inference.} The visual design aims for a restrained instrument-panel aesthetic rather than a decorative smart display aesthetic. The choice of typography, centered composition, and sparse iconography supports that reading.

\section{Power and Sleep Architecture}
\subsection{Setup-Driven Runtime Model}
The firmware is effectively \texttt{setup()}-driven. The device wakes, restores state, performs any required work, draws the screen, and returns to deep sleep. The \texttt{loop()} function is not used for continuous normal operation and simply routes back into sleep.

\subsection{Wake Timing}
The device wakes at approximately 58 seconds past each minute. This creates a short pre-minute execution window in which the firmware can:
\begin{itemize}
  \item check for button hold requests,
  \item perform an hourly weather update if needed,
  \item defer display initialization until the appropriate render moment.
\end{itemize}

This scheme allows the visible display update to remain aligned with the minute boundary while still giving the processor enough time to prepare.

\subsection{Full Versus Partial Refresh}
The refresh strategy is:
\begin{itemize}
  \item full refresh on cold boot,
  \item full refresh at the top of each hour,
  \item partial time-only refresh on other minute transitions.
\end{itemize}

The time refresh region is precomputed and aligned to 8-pixel boundaries, which matches the partial window constraints of the display controller.

\subsection{Display Initialization Deferral}
One of the most technically significant design choices is that display initialization is deferred during timer wakes. The repository documents that early initialization would pulse the display reset line and cause a visible artifact at the \texttt{:58} wake moment. Deferring \texttt{display.init()} until just before drawing prevents an unnecessary visible flash and makes the wake cycle effectively silent from the user perspective.

\subsection{GPIO Hold Strategy}
Before deep sleep, the firmware enables GPIO hold on the display control lines associated with reset and signaling. This prevents the ESP32-C6 boot ROM from briefly glitching those lines during the next wake sequence, which would otherwise wake the e-paper controller and create a visible flash.

This is a thoughtful hardware-software integration detail and represents one of the strongest technical aspects of the project.

\section{State Management}
\subsection{Persistent Configuration}
Long-lived user configuration is stored in ESP32 NVS under a dedicated namespace. The stored settings include:
\begin{itemize}
  \item Wi-Fi SSID,
  \item Wi-Fi password,
  \item POSIX timezone string,
  \item corresponding IANA timezone name,
  \item latitude,
  \item longitude,
  \item dark mode flag.
\end{itemize}

\subsection{RTC Memory State}
Transient but sleep-persistent runtime state is stored in RTC memory using \texttt{RTC\_DATA\_ATTR}. This includes:
\begin{itemize}
  \item weather validity and most recently fetched values,
  \item last successful weather fetch timestamp,
  \item dark mode working state,
  \item cached city name,
  \item validity flag for the RTC-backed state block.
\end{itemize}

This architecture reduces network activity and improves resilience across minute wake cycles.

\section{Networked Features}
\subsection{NTP Synchronization}
On cold boot the device performs NTP synchronization using public time servers. Before SNTP is started, the firmware explicitly resets the ESP32 system clock to zero. The repository notes that without doing so, stale system time retained in RTC slow memory could be misinterpreted as a fresh NTP result.

This is a subtle but important correctness measure.

\subsection{Weather API}
Weather data is requested from Open-Meteo over plain HTTP. The project notes explain that HTTPS was avoided for this specific endpoint because the ESP32-C6 TLS stack exhibited handshake reliability issues against the service configuration, while the weather data is public and unauthenticated.

From a product perspective, this is a pragmatic engineering tradeoff:
\begin{itemize}
  \item reliability of fetches was prioritized,
  \item no secrets are exposed by the request itself,
  \item the endpoint provides non-sensitive public data.
\end{itemize}

\subsection{Geocoding}
Reverse geocoding and setup-portal city search use Nominatim over HTTPS, with a custom user agent string and relaxed certificate verification through \texttt{setInsecure()}.

\textbf{Labeled inference.} For a personal desk device, this is a reasonable integration choice, though a more security-constrained product version would likely replace permissive TLS verification with pinned or properly validated certificates.

\section{Configuration Portal}
\subsection{Entry Method}
The setup portal is entered by holding the button for approximately one second while the processor is awake. The current design does not support true deep-sleep button wake on the chosen input pin because GPIO 17 is not a low-power wake-capable pin on the ESP32-C6.

\subsection{Portal Behavior}
When activated, the firmware:
\begin{itemize}
  \item draws a setup screen on the e-paper display,
  \item starts a local Wi-Fi access point,
  \item enables a DNS catch-all responder for captive portal behavior,
  \item serves a web configuration interface on port 80,
  \item optionally connects to the previously saved home network in STA mode for city search proxying.
\end{itemize}

The portal times out after five minutes if no configuration is saved.

\subsection{User-Facing Setup Functions}
The web UI allows the user to:
\begin{itemize}
  \item confirm Wi-Fi connectivity state,
  \item enter or preserve Wi-Fi credentials,
  \item search for a city and store latitude/longitude,
  \item choose a timezone from a curated list,
  \item toggle dark mode,
  \item save and restart the device.
\end{itemize}

\section{Boot and Execution Flow}
\subsection{Cold Boot Path}
The cold boot path can be summarized as:
\begin{enumerate}
  \item load saved configuration,
  \item initialize hardware interfaces,
  \item initialize the display immediately,
  \item initialize the RTC,
  \item check for a button-held portal request,
  \item compute layout geometry,
  \item connect to Wi-Fi,
  \item synchronize time from NTP to the RTC,
  \item fetch city name,
  \item fetch weather,
  \item render a full screen,
  \item enter deep sleep.
\end{enumerate}

\subsection{Minute Wake Path}
Normal minute wake operation can be summarized as:
\begin{enumerate}
  \item wake at \texttt{:58},
  \item load configuration and restore sleep-persistent state,
  \item release held GPIO lines,
  \item check whether setup mode is requested,
  \item if appropriate, fetch hourly weather update,
  \item wait for the safe draw window,
  \item initialize the display only when draw time is reached,
  \item perform either full or partial refresh,
  \item return to deep sleep.
\end{enumerate}

\section{Current Limitations}
The current implementation is already functional and coherent, but it includes several known limitations:
\begin{itemize}
  \item button entry to setup mode is only available during the active wake window,
  \item battery level is not yet measured or displayed in the production firmware,
  \item weather refresh cadence is fixed rather than adaptive,
  \item the portal includes no image assets, diagrams, or richer diagnostics,
  \item networking is intentionally opportunistic rather than continuously connected.
\end{itemize}

\section{Product Characterization}
If treated as a small product rather than only a hobby prototype, the current firmware can be characterized as follows.

\subsection{Strengths}
\begin{itemize}
  \item Strong low-power architecture based on deep sleep and RTC ownership of time.
  \item Careful e-paper behavior management, especially around reset artifacts.
  \item User-configurable setup without requiring firmware recompilation.
  \item Clean information hierarchy on screen.
  \item Limited but purposeful network dependency.
\end{itemize}

\subsection{Risks for Productization}
\begin{itemize}
  \item Field update strategy is not yet present.
  \item Battery behavior is not yet fully integrated.
  \item Some external services are public community services and may not be suitable unchanged for scaled deployment.
  \item TLS verification is permissive for geocoding requests.
\end{itemize}

\section{Build and Development Notes}
The repository indicates a standard PlatformIO workflow for local builds and flashing. Development currently appears optimized for direct firmware iteration rather than for a release pipeline with automated validation.

\textbf{Labeled inference.} This is appropriate for the project stage. The codebase reads like a carefully iterated personal device firmware rather than a generalized SDK or mass-production codebase.

\section{Planned and Future Work}
This section records features that are either explicitly staged in the repository or recommended as natural next steps for the product.

\subsection{Planned Battery Measurement Integration}
The repository includes a non-integrated battery feature draft in \texttt{battery-measuring/battery\_code\_backup.cpp}. The intended design uses a resistor divider wired to an ADC input so the firmware can estimate pack voltage and derive battery percentage.

The staged implementation includes:
\begin{itemize}
  \item voltage sampling through the ADC,
  \item a discharge-curve-based percentage estimate,
  \item charging-state detection,
  \item an on-screen battery icon or charging indicator,
  \item a dedicated low-battery or dead-battery screen,
  \item extended sleep behavior when the battery is critically low.
\end{itemize}

This is the most obvious near-term enhancement because much of the logic has already been conceptually specified and partially written. Once the resistor divider is physically wired as required, the feature should integrate naturally into the existing hourly full-refresh flow and low-power architecture.

\subsection{Recommended Adaptive Night Refresh Policy}
An additional recommended feature is reduced refresh frequency during nighttime hours. The existing firmware updates the time every minute, but e-paper clocks do not necessarily need identical visual update cadence across the full 24-hour day.

A future adaptive policy could:
\begin{itemize}
  \item retain the current minute refresh rate during daytime,
  \item switch to a slower cadence during configured night hours,
  \item continue to preserve hourly weather maintenance if desired,
  \item reduce battery consumption and unnecessary display activity.
\end{itemize}

This feature aligns strongly with the current design philosophy because it improves power efficiency without changing the visual identity or core use model of the device.

\subsection{Other Viable Future Enhancements}
Other future additions that would fit the product well include:
\begin{itemize}
  \item alternate visual themes or layout modes,
  \item indoor environmental sensing such as room temperature,
  \item OTA update support,
  \item sunrise, sunset, or moon-phase indicators,
  \item lightweight calendar or reminder overlays.
\end{itemize}

\section{Conclusion}
\texttt{shravclock} is a thoughtfully engineered personal desk clock with a notably strong low-power and e-paper-oriented architecture. Its most mature qualities are the wake-refresh-sleep execution model, the deliberate handling of display hardware quirks, the RTC-centered time strategy, and the inclusion of a self-contained configuration portal.

At its current stage, the project already operates as a complete and polished personal device. The most natural next milestone is battery instrumentation, followed by adaptive nighttime behavior and broader product-hardening improvements. As documented in this report, the project is already substantially beyond a simple clock build; it is a compact embedded product with clear architecture, deliberate tradeoffs, and a credible path toward a more fully featured standalone object.

\end{document}
